%% §1 — the problem, the sentence, the design in brief, the kernel,
%% the measurements, contributions, and what is not claimed. Written
%% against tag era6-campaign-3; vocabulary is KERNEL.md §12's.

\section{Introduction}
\label{sec:intro}

Verification infrastructure is now being written by language models:
interpreters for the formats benchmarks arrive in, translators between
them, invariant generators, proof scripts, and the checkers that are
supposed to catch the rest. Each such artifact is text produced by a
process that cannot be inspected and that is wrong at an unknown rate,
and the question every deployment faces is what a result computed by
such artifacts is worth. Two answers are current. One verifies the
artifact: repair the proof until a checker accepts it, or generate
the code together with a specification and a proof
\citep{first2023baldur,sun2024clover}. The other distrusts the
artifact and checks its outputs: certifying algorithms
\citep{mcconnell2011certifying}, proof-carrying code
\citep{necula1997pcc}, validated verification witnesses
\citep{beyer2015witnesses}. Both answers are about one artifact at a
time. Neither says what happens when the checker is generated too,
when the translation that carried the program to the checker is
generated, when the map that carries the checker's verdict back to
the original program is generated, and when the search that produced
the evidence was generated last week by the same process.

This paper describes a system built on the premise that every one of
those pieces is untrusted, and judged whole --- and built that way in
fact, not as a thought experiment. It started from an empty
repository with a specification and nothing else, and everything that
has entered since --- every interpreter, translator, search, checker,
and carry-back --- is Python written by a language model and admitted
through one gate. No human wrote any of it, and no human vouches for
it line by line; what stands in for that is the subject of the paper.
The separation that organizes it is a single sentence:

\begin{quote}
\textbf{Generation produces syntax; only interpretation produces
truth.}
\end{quote}

A \emph{language} is a syntax with an \emph{interpreter} --- the
program that runs its programs and reports named observables --- and,
where the language declares \emph{certificate forms}, a
\emph{checker} for each. Interpreters and checkers are the
\emph{judges}, and the trusted base of the system is exactly the set
of admitted judges, a list the kernel prints. Every other executable
is a \emph{transport}: a \emph{translator} from one language to
another, a \emph{carry-back} that brings evidence found at the target
home to the source, or a \emph{search} that looks for evidence --- a
\emph{witness} or a \emph{certificate} --- within a budget. A
transport is untrusted whatever it computes; its output is only ever
as good as the judgment it survives. \textbf{Searches do not decide;
they write. Judges decide.}

The consequence for a result is that its \emph{grade} is geometry. A
question lives at a language; a \emph{route} carries it across
\emph{pairs} of languages to a search; whatever comes back is judged
where it lands, and the \emph{gap} of a result is the number of hops
between the question and the last judgment its evidence passed. Gap
zero is \grade{certified}: a witness replayed by the interpreter of
the language the question lives in, or a certificate discharged by
that language's checker. A positive gap is \grade{checked}: the
certificate was discharged one or more hops away, and the result
still rests on the translators between. No judgment at all is
\grade{claimed}. The \emph{residual} of a result --- what it still
rests on --- is computed from the log as the lineages of the hops
inside the gap plus the judge that ran; every judgment removes
everything upstream of it, so at gap zero the search that found the
evidence, and the route it travelled, are not in the residual at all
(\cref{thm:residual}, proved in Lean).

A design like this makes two claims that can be wrong. The first is
about transports: a faulty translator, carry-back, or search can cost
answers but can never forge a \grade{certified} one, and whatever it
does corrupt says so in its grade and names it in its residual. The
second is about judges: the trusted base is small enough, and tested
well enough, to be believed. We test both by fault injection
(\cref{sec:faults}). For the first we mutate every transport in the
registry, \emph{remove the gate}, force each mutant in, and play it
on questions whose answer is known: 16,720 plays of 1,672 mutants
yield 133 wrong answers, 127 of them \grade{claimed}, 6
\grade{checked} at gap one, and none \grade{certified}; all 133 carry
the faulty entry's lineage in their residual, and seven of the
mutants that answered wrongly had passed the gate. For the second we
mutate the judges and let the gate, and everything admitted
downstream of a judge, try to refuse them: half of 700 are refused,
the survivors concentrate in the certificate checkers, and a
differential probe finds six of them laxer than the checker they
mutate --- real holes in the trusted base. We close the holes without
touching a judge, by revisions that add test vectors whose expected
values are confirmed by tools that never enter the system: a model
checker, a compiler, and the formal model of an instruction set.

The system's own boards are the third measurement
(\cref{sec:boards}). Two domains have entered: hardware, whose root
language is BTOR2 \citep{niemetz2018btor2}, and software, whose root
is a fragment of C; RISC-V is a third language, anchored by no domain
and corroborated by the pairs on either side of it. Four generated
pairs connect the three languages, six generated searches live at
BTOR2 and C, and four benchmarks pinned from HWMCC'24 and SV-COMP
2025 stand at 45 of 74, 26 of 79, 47 of 80, and 12 of 55 questions
settled, with 49 universal answers \grade{certified} at gap zero and
no contradiction. The generated searches are not competitive with the
solvers whose verdicts sit beside the boards, a translation of
hardware into C returned a partial on every route it opened, and ten
times the clock bought one question; the paper reports these because
the system cannot say otherwise.

\paragraph{Provenance}
The premise applies to this paper as well. All of the system --- the
kernel, its tests, the second lineage of its pure half, the Lean
development, every entry in the registry --- and the instruments of
\cref{sec:measurements} --- the fault-injection harness, the
generators of test-vector expectations --- are code generated by
language-model agents (Anthropic's Claude Fable, through the Claude
Code command-line agent). Most of this paper's text, its tables, and
its figures are generated by the same agents, from the specification
and the measurement records. No human reviewed the code or the text
for semantic correctness. The human contribution is the
architecture: the separation of judges from transports, the gate, the
decision that grades are computed and never asserted, and the choices
of what to build, what to measure, and what to claim. What the reader
is offered in place of trust in the writers is what the system offers
in place of trust in its transports. No number here originates with
a model: each is read from logs and stamps the kernel wrote, the
boards regenerate from those logs byte for byte, the fault-injection
counts are recomputed from the deposited records by one command, the
proofs are checked by Lean and their axiom audit is printed at every
build, and the test vectors' expected values were confirmed by tools
that are not generative. The reader need not trust
the writers, human or otherwise, beyond what the thesis itself
requires.

\paragraph{Contributions}
\begin{itemize}
\item A design in which every executable of a verification system is
  untrusted generated text and one gate admits all of them: two
  kinds (language, transport), three judgments an artifact faces on
  crossing a pair (the square, replay, discharge), and two crossings
  nothing judges (\cref{sec:design}).
\item Grades as geometry: gap, grade, and residual computed from the
  log, with the order on results and the residual's properties proved
  (\cref{sec:grades,sec:kernel}).
\item A kernel outside the gate made solid four ways --- proof,
  tests, a clean-room second lineage, content pins --- each worded no
  stronger than what it verifies (\cref{sec:kernel}).
\item A fault-injection method for such systems, and its results: the
  no-forge claim tested with the gate removed, the trusted base's
  exposure measured and localized, six holes found and closed by
  vector-only revisions with outside ground truth (\cref{sec:faults}).
\item Measurements at one frozen commit, every number tied to a file
  in the tree, the negatives reported as results
  (\cref{sec:boards}).
\end{itemize}

\paragraph{What is not claimed}
The square is checked empirically, per program, on a pair's corpus;
per-program translation correctness is not proved. The seal around
generated executables is a blank \code{PATH} and an emptied
environment, measured by tests, and not claimed to be a proof. Half
the sampled judge mutants survive admission; for most of them we can
say only that no certificate we could find distinguishes them from
the judge. The \emph{corroborated} flag has fired only between two
searches, never between two judges. And the boards are not compared
to any solver's on a shared clock; the competitions' verdicts are
recorded beside each question as testimony, never ranked.
